% TOP_PROBLEM: {"id": 4900003, "title": "Banach's simple Lebesgue spectrum problem", "category_id": 11, "category": {"id": 11, "name": "dynamical_systems", "display_name": "Dynamical Systems", "description": "Problems about long-term behavior of deterministic systems, Hamiltonian dynamics, and periodic orbits.", "slug": "dynamical-systems", "order_index": 11, "created_at": "2026-07-31T15:26:25.671Z"}, "status": "partially_solved"}
% FIRST_POSTED: 2026-10-06
% ATTEMPT_STATUS: solved
% OPENAI_CLAIM: {"schema_version":1,"source_repo":"https://github.com/openai/math","source_commit":"adc7f1241b42e322a6451854ab7e4b4c146bf78a","release_date":"2026-10-06","match":"stronger","resolution":"proved","independently_reviewed":false,"second_pass":"agreed","families":[{"family":"144","title":"Banach’s simple Lebesgue-spectrum problem","subject":"Dynamical systems and ergodic theory","manuscripts":[{"title":"A smooth three-torus diffeomorphism with simple Lebesgue spectrum","pdf_path":"preprints/A-smooth-three-torus-diffeomorphism-with-simple-Lebesgue-spectrum-September-23-2026/paper.pdf","readme_path":"preprints/A-smooth-three-torus-diffeomorphism-with-simple-Lebesgue-spectrum-September-23-2026/README.md","date":"2026-09-23","theorem_ref":"Theorem labelled thm:main"}],"lean":{"doc_path":"lean/docs/144.md","comparators":["lean/ComparatorChallenges/ThreeTorus.json","lean/ComparatorChallenges/ThreeTorus.lean"],"declarations":[],"covers_main_theorem":true},"reasoning_trace":null}]}
\documentclass[11pt]{article}
\usepackage[margin=25mm]{geometry}
\usepackage{fontspec}
\setmainfont{Latin Modern Roman}
\usepackage{amsmath,amssymb,mathtools}
\usepackage{unicode-math}
\setmathfont{Latin Modern Math}
\usepackage{enumitem}
\usepackage{xurl,hyperref,bookmark}
\hypersetup{colorlinks=true,urlcolor=blue}
\setcounter{secnumdepth}{0}
\setlength{\parindent}{0pt}
\setlength{\parskip}{5pt}
\setlength{\emergencystretch}{3em}
\begin{document}
\section*{Banach's simple Lebesgue spectrum problem}

\subsection{OpenAI's claim}
A smooth volume-preserving diffeomorphism of the three-torus is claimed to have simple Lebesgue spectrum on its entire mean-zero square-integrable space. The cited manuscript argues for that result; its relation to this problem is specified below.

\subsection{Outline of the argument}
Manuscript 1 (family 144, 2026-09-23): Finite blocks of small density-weighted biases learn observable labels while losing only a quadratic amount of spectral mass. Whitening restores density one with a small physical-space change; a diagonal construction preserves an orthonormal orbit and makes that orbit span a dense family of observables.

\subsection{Scope relative to this problem}
Manuscript 1 (family 144, 2026-09-23): The smooth volume-preserving three-torus example has a bilateral orthonormal cyclic basis on the full mean-zero L2 space, satisfying Banach’s usual simple-Lebesgue-spectrum formulation with additional smooth realization.

\subsection{Formal verification}
For family 144, the scope page documents the following selections. The torus construction is selected on the full mean-zero square-integrable space, including a real cyclic vector with an orthonormal integer orbit and the circle-multiplication model. The documented statement covers the main claim of manuscript 1. Its selected declaration also covers the clause used for this problem. Comparator files: ThreeTorus.lean. Declarations: OAI.ThreeTorus.simple\_lebesgue\_spectrum.

This site has not executed the Lean files or checked their dependencies. These descriptions report the scope of the linked declarations, rather than a new verification of the proofs.

\subsection{Status caveat}
OpenAI attributes these results to an internal model. Its release notice cautions that results lacking full formalization can contain mistakes. This site has not independently assessed the proofs; a solved label records the scope of the reported claim.

\subsection{Sources}
\begin{itemize}
\item OpenAI, \emph{A smooth three-torus diffeomorphism with simple Lebesgue spectrum}, 2026-09-23; citation key \texttt{OAI:A-smooth-three-torus-diffeomorphism-with-simple-Lebesgue-spectrum-September-23-2026}. \href{https://github.com/openai/math/blob/adc7f1241b42e322a6451854ab7e4b4c146bf78a/preprints/A-smooth-three-torus-diffeomorphism-with-simple-Lebesgue-spectrum-September-23-2026/paper.pdf}{PDF (GitHub)}; \href{https://github.com/openai/math/raw/adc7f1241b42e322a6451854ab7e4b4c146bf78a/preprints/A-smooth-three-torus-diffeomorphism-with-simple-Lebesgue-spectrum-September-23-2026/paper.pdf}{PDF (direct)}; \href{https://github.com/openai/math/blob/adc7f1241b42e322a6451854ab7e4b4c146bf78a/preprints/A-smooth-three-torus-diffeomorphism-with-simple-Lebesgue-spectrum-September-23-2026/README.md}{Manuscript page and citation}; \href{https://github.com/openai/math/blob/main/preprints/A-smooth-three-torus-diffeomorphism-with-simple-Lebesgue-spectrum-September-23-2026/paper.pdf}{Latest version}.
\item Family 144: \href{https://github.com/openai/math/blob/adc7f1241b42e322a6451854ab7e4b4c146bf78a/lean/docs/144.md}{Lean scope}.
\item \href{https://github.com/openai/math/blob/adc7f1241b42e322a6451854ab7e4b4c146bf78a/lean/ComparatorChallenges/ThreeTorus.json}{Comparator ThreeTorus.json}.
\item \href{https://github.com/openai/math/blob/adc7f1241b42e322a6451854ab7e4b4c146bf78a/lean/ComparatorChallenges/ThreeTorus.lean}{Comparator ThreeTorus.lean}.
\item \href{https://github.com/openai/math/blob/adc7f1241b42e322a6451854ab7e4b4c146bf78a/CONTENTS.md}{OpenAI catalogue entry}.
\end{itemize}
\end{document}
